\section{中央 pencil の例外代数と圏論的 blow-up の構成}
\label{sec:restored-rees-sod}
前節の Rees 環の表示からは、例外因子の代数と、
\emph{比較対象としての} categorical blow-up をさらに構成できる。
しかし両者を同一視するには、別個の比較定理が必要である。

\begin{proposition}[随伴次数環]\label{prop:restored-associated-graded}
第\ref{prop:rees-central-regular}命題の仮定の下で $D=A/(s,t)$ とおく。
Rees 次数を $m$ とすると、自然な同型
\[
\operatorname{gr}_{I}A=\bigoplus_{m\ge0}I^m/I^{m+1}
 \simeq D[U,V],\qquad [s]\mapsto U,\ [t]\mapsto V
\]
がある。右辺で $U,V$ は $D$ と可換な中央不定元である。
特に $D$ が次数六の楕円 twisted homogeneous coordinate ring で、
$\operatorname{qgr}D\simeq\operatorname{coh}E$ なら、
Rees 方向の例外射影模型の局所チャートは
$D[w]$ と $D[w^{-1}]$ で記述できる。
偏極の共通 twist を除くと、例外モデルは
$E\times\mathbb P^1$ の非可換代数的表示になる。
\end{proposition}
\begin{proof}
Rees 代数 $\mathcal R_A(I)$ を、基底環 $A$ から来る
次数零の二側イデアル $I\mathcal R_A(I)$ で割ると
\[
\mathcal R_A(I)/I\mathcal R_A(I)
 =\bigoplus_{m\ge0}I^m/I^{m+1}
\]
を得る。
第\ref{prop:rees-central-regular}命題の表示ではこの商は
\[
 A[U,V]/(tU-sV,s,t)=(A/I)[U,V]
\]
である。$U$ と $V$ は中央なので、$U\ne0$ と $V\ne0$ の
二つの標準チャートを取ると、それぞれの相対次数零部分は
$D[V/U]$ と $D[U/V]$ になる。
交叉上で両座標が互いに逆になるから、最後の主張が従う。
\end{proof}

\begin{theorem}[制限関手からの categorical blow-up]\label{thm:restored-categorical-gluing}
$\mathcal C$ と $\mathcal D$ を smooth proper な冪等完備 dg 圏とし、
$\operatorname{Res}:\mathcal C\to\mathcal D$ を perfect dg bimodule によって
表される dg 関手とする。
すると smooth proper な dg 圏 $\mathcal B_{\mathrm{cat}}$ が構成でき、
\[
 \mathcal B_{\mathrm{cat}}
 =\langle\mathcal D,\mathcal C\rangle
\]
という admissible な半直交分解を持つ。
ここで混合 Hom は、$G\in\mathcal D,F\in\mathcal C$ に対して
\[
\operatorname{RHom}_{\mathcal B_{\mathrm{cat}}}(G,F)
 \simeq\operatorname{RHom}_{\mathcal D}(G,\operatorname{Res}F)[-1],
\qquad
\operatorname{RHom}_{\mathcal B_{\mathrm{cat}}}(F,G)=0
\]
と指定できる。
さらに $X$ を滑らかな射影多様体、$Z\subset X$ を滑らかな余次元二の中心とし、
$\mathcal C=\operatorname{Perf}(X)$、$\mathcal D=\operatorname{Perf}(Z)$、
$\operatorname{Res}=Li^*$ とすると、
この圏は標準的な dg enhancement の意味で
$\operatorname{Perf}(\operatorname{Bl}_Z X)$ と同値である。
\end{theorem}
\begin{proof}
まず $\mathcal D$ と $\mathcal C$ を二つの対角成分とし、
$\mathcal D$ から $\mathcal C$ への Hom を指定の perfect bimodule、
逆向きの Hom を零とする上三角 dg category を作る。
その perfect closure を $\mathcal B_{\mathrm{cat}}$ とおく。
上三角性から二つの fully faithful な包含と表示した半直交分解が得られる。
両成分の smooth properness と bimodule の perfect 性から、
Orlov の gluing 定理\cite{PencilOrlovGluing}によって
$\mathcal B_{\mathrm{cat}}$ も smooth proper である。

可換の場合の同定を確認する。
$f:Y=\operatorname{Bl}_Z X\to X$、$j:P(N_{Z/X})\hookrightarrow Y$、
$p:P(N_{Z/X})\to Z$ とおく。
$\Psi(G)=j_*(p^*G\otimes\mathcal O_{P(N)}(-1))$ とすれば、
Orlov の blow-up 公式は
$\operatorname{Perf}(Y)=\langle\Psi\operatorname{Perf}(Z),
Lf^*\operatorname{Perf}(X)\rangle$ を与える\cite{OrlovBlowup}。
$Rp_*\mathcal O_{P(N)}(-1)=0$ により
$\operatorname{RHom}_Y(f^*F,\Psi G)=0$ である。
また、例外因子の正規束が $\mathcal O_{P(N)}(-1)$ なので
$j^!f^*F\simeq p^*i^*F\otimes\mathcal O_{P(N)}(-1)[-1]$。
従って
\[
\operatorname{RHom}_Y(\Psi G,f^*F)
 \simeq\operatorname{RHom}_{Z}(G,Li^*F)[-1].
\]
この dg bimodule と両対角成分が一致するから、
gluing の普遍性により主張の同値が従う。
\end{proof}

\begin{remark}[中央 pencil への適用範囲]\label{rem:restored-sod-scope}
最大 pencil の代数 $A$ で $A/(s,t)=B(E,M,\tau)$ とすると、
$E$ は中心の候補である。
$\mathcal C=\operatorname{Perf}(\operatorname{qgr}A)$ が smooth proper で、
中央商に沿う導来制限 $\operatorname{Res}:\mathcal C\to\operatorname{Perf}(E)$ が
perfect bimodule で表されるなら、上の定理は
$\langle\operatorname{Perf}(E),\mathcal C\rangle$ を持つ
\emph{categorical} blow-up の存在を与える。
しかしこの圏が第\ref{prop:rees-central-regular}命題の
Rees 環の $\operatorname{qgr}$ から作る圏と同値であることは、
本書ではまだ証明していない。
また $\operatorname{Perf}(\operatorname{qgr}A)$ と
第\ref{pcat:sec}節の階数64の mutation 代数の同値も未解決である。
これらを混同すると、存在が未確認の SOD を Rees 模型へ誤って付与することになる。
\end{remark}

\begin{proposition}[generic localization の条件付き両立性]
上の $\mathcal B_{\mathrm{cat}}=\langle\mathcal D,\mathcal C\rangle$ において、
$L:\mathcal C\to\mathcal G$ を Verdier 局所化とする。
$\mathcal B_{\mathrm{cat}}/\mathcal D$ は冪等完備化を除き $\mathcal C$ と同値であり、
従って $L$ は $\mathcal B_{\mathrm{cat}}$ から $\mathcal G$ への
局所化へ合成され、$\mathcal D$ の対象はそこで零になる。
ただしこの事実だけから、別の Rees 圏の generic localization と同一視することはできない。
\end{proposition}
\begin{proof}
admissible SOD の Verdier 商の普遍性による。
\end{proof}
